\documentclass[11pt]{article}

\usepackage[T1]{fontenc}
\usepackage[utf8]{inputenc}
\usepackage[margin=1in]{geometry}
\usepackage{microtype}
\usepackage{titling}

\setlength{\droptitle}{-3.5em}

\title{\vspace{-1em}QueryBloom: Automated Text--SPARQL Example Synthesis\\ via Subgraph Sampling}

\author{
  Essam Abdelghany and Chengkai Li \\[0.3em]
  \normalsize Department of Computer Science and Engineering \\
  \normalsize The University of Texas at Arlington \\
  \normalsize Field of Study: Computer Science (Natural Language Processing / Machine Learning)
}
\date{}

\begin{document}
\maketitle
\vspace{-2.5em}

Knowledge graphs (KGs) store vast amounts of structured, domain-specific knowledge, and question answering systems that generate SPARQL queries from natural language empower researchers to directly access such knowledge. Training, tuning, and
evaluating such systems depends on Text--SPARQL example sets, which supply the
in-context demonstrations, fine-tuning data, and evaluation benchmarks these systems
require. Building these sets by hand is
slow and costly, since it requires annotators fluent in both the target domain and the
KG's ontology. As a result, curated examples exist almost exclusively for a few
open-domain graphs, leaving domain-specific KGs in various fields like biomedicine and
geoscience with little coverage.

We address this gap with \textbf{QueryBloom}, an automated, KG-agnostic method that
synthesizes diverse, execution-verified Text--SPARQL examples from a given RDF knowledge
graph using only its SPARQL endpoint. QueryBloom builds on the observation that every
SPARQL query implicitly describes a subgraph pattern over the graph, and inverts this
correspondence. It first mines a diverse set of type patterns that capture which properties connect which kinds of entities, then grows structurally diverse multi-hop subgraphs by extending these patterns, verifying against the endpoint that each retains sufficient support. A reflective
large language model then verbalizes each grounded subgraph pattern into a natural language
question and a corresponding SPARQL query, using execution feedback to iteratively
verify and correct every pair.


We evaluated QueryBloom against the human-authored DBpedia QALD-9 and
machine-generated DBLP-QuAD benchmarks, finding that it matches or exceeds both in
query complexity, linguistic variety, and schema coverage. Because each example
combines repeated endpoint verification with iterative model inference, the total computational
cost scales with the target dataset size, making high-performance resources such as those at TACC imperative for practical dataset generation at scale. By requiring only a query endpoint, QueryBloom makes it practical to
bootstrap high-quality KGQA datasets for domain-specific knowledge graphs that
previously had none.

\end{document}
